\documentclass[a4paper,10pt]{article} \usepackage{anysize}
\marginsize{2cm}{2cm}{1cm}{1cm}
%\textwidth 6.0in \textheight = 664pt
\usepackage{xltxtra}
\usepackage{xunicode} 
\usepackage{graphicx}
\usepackage{color} 
\usepackage{xgreek} 
\usepackage{fancyvrb}
\usepackage{minted}
\usepackage{listings}
\usepackage{enumitem} 
\usepackage{framed} 
\usepackage{relsize}
\usepackage{float} 
\usepackage{fontspec}
\setmainfont[Mapping=TeX-text]{DejaVu Serif}

%\setmainfont{Arial}
%\setsansfont{FreeSans}
%\setmonofont{FreeMono}
\begin{document}
%\def\thesubsection {\alph{subsection}}
\renewcommand{\labelenumi}{\alph{enumi}}
\renewcommand{\labelenumii}{(\roman{enumii})}

\begin{titlepage}
\begin{center}
\begin{figure}[t] 
     \includegraphics[scale=0.7]{title/ntua_logo}
\end{figure}
\begin{LARGE}\textbf{ΕΘΝΙΚΟ ΜΕΤΣΟΒΙΟ ΠΟΛΥΤΕΧΝΕΙΟ\\}\end{LARGE}
\vspace{5cm}
\begin{Large}
ΣΧΟΛΗ ΗΜ\&ΜΥ\\
Αλγόριθμοι και Πολυπλοκότητα\\
1\textsuperscript{η} Σειρά Γραπτών Ασκήσεων\\
Ακ. έτος 2010-2011\\
\end{Large}
\vspace{10cm}
\begin{tabular}{l r}
\Large{Λύρας Γρηγόρης}&
\large{Α.Μ.: 03109687}\\
\end{tabular}\\
\vspace{1cm}

\vfill
\large\today\\
\end{center}
\end{titlepage}



\section{Ασυμπτωτικός συμβολισμός, Αναδρομικές
Σχέσεις} \setcounter{section}{1}
\begin{enumerate}
\item Ταξινόμηση
\begin{tabular}{l l}
$ n*3^n$ & $\Rightarrow O(3^n)$ \\
$ n^{1.01}$ & $\Rightarrow O(n^{1.01})$\\
$ 5^{\log _2 {n}}$ & $\Rightarrow O(n^{2.321})$\\
$ \sum_{k=1}^{n}k^5$ & $\Rightarrow O(n^6)$\\
$ 2^{{\log _2 {n}}^4}$ & $\Rightarrow O(n^4)$\\
$ {\log^{\log{n}}{n}}$ & $\Rightarrow O()$\\
$ \frac{n}{\log\log{n}}$ & $\Rightarrow O()$\\
$ \exp{\frac{n}{\ln{n}}}$ & $\Rightarrow O()$\\
$ \log{n^3}$ & $\Rightarrow O()$\\
$ \sqrt{n}*(\log{n})^{50}$ & $\Rightarrow O()$\\
$ n*(\log{n})^{10}$ & $\Rightarrow O()$\\
$ (\log{n})^{\sqrt{n}}$ & $\Rightarrow O()$\\
$ n^{\log{\log{n}}}$ & $\Rightarrow O()$\\
$ 2^{2*n}$ & $\Rightarrow O()$\\
$ \sqrt{n!}$ & $\Rightarrow O()$\\
$ \log{(n!)}$ & $\Rightarrow O(n*\log{n})$\\

\end{tabular}

\item Τάξη Μεγέθους
\begin{enumerate}
\item $T(n)=5*T(n/7)+n*\log(n)$
$\Rightarrow n^{\log _5 {7}} = n^{1.209}$
$\Rightarrow n^{1.209} > n*\log{n}$\\
$\Rightarrow T(n) \in \Theta(n^{1.209})$

\item $T(n) = 4*T(n/5)+n/\log^2{n}$
$\Rightarrow n^{\log_4 {5}} = n^{1.16}$
$\Rightarrow n^{1.16} > n/\log^2{n}$\\
$\Rightarrow T(n) \in \Theta(n^{1.16})$

\item $T(n) = T(n/3) +3*T(n/7)+n$\\
$\Rightarrow T(n) \in O(n*k^{\log{n}}*\log{n})$

\item $T(n) = 6*T(n/6)+n$
$\Rightarrow n^{\log _6 {6}} = n$\\
$\Rightarrow T(n) \in \Theta(n*\log _6 {n})$

\item $T(n) = T(n/3) + T(2n/3) + n$\\
$\Rightarrow T(n) \in \Theta(n*\log{n})$

\item $T(n) = 16*T(n/4) + n^3*\log^2{n}$
$\Rightarrow n^{\log _16{4}} = \sqrt{n}$
$\Rightarrow \sqrt{n} < n^3*\log^2{n}$\\
$\Rightarrow T(n) \in \Theta(n^3*\log^2{n})$

\item $T(n) = T(\sqrt(n) + \Theta(\log{\log{n}})$

\item $T(n) = T(n-3) + \log{n}$\\
$\Rightarrow T(n) \in \Theta(n*logn)$

\end{enumerate}
\end{enumerate}

\section{Ταξινόμιση σε Πίνακα με Πολλά Ίδια Στοιχεία}

\section{Δυαδική Αναζήτηση}
\begin{enumerate}
\item Μπορούμε να σπάμε τον πίνακα σε $n/2$ στοιχεία σε κάθε στάδιο. Έπειτα να
ελέγχουμε το στοιχείο $n/2$ και το στοιχείο $n/2+1$ με το $x$. Αν είναι
αν είναι μικρότερο του μικρότερου επαναλαμβάνουμε στο πρώτο τμήμα. Αν είναι
μεγαλύτερο του μεγαλύτερου επαναλαμβάνουμε στο δεύτερο τμήμα. Αν δεν ισχύει
τίποτα από τα δύο, δεν υπάρχει θέση που να περιέχει τον $x$.
\item Παίρνω τα k πρώτα στοιχεία από κάθε πίνακα. Τα χωρίζω στη μέση και
συγκρίνω τα $A[k/2]$ με $B[k/2+1]$ και $A[k/2+1]$ με $B[k/2]$. Αν
$A[k/2]<B[k/2+1]$ και $B[k/2]<A[k/2+1]$ τοτε το k-οστό στοιχείο είναι το
$max\{A[k/2],B[k/2]\}$. Αν $A[k/2]<B[k/2+1]$ και $B[k/2]>A[k/2+1]$ τότε
επαναλαμβάνω χρησιμοποιώντας τα στοιχεία $[A[k/2+1]..A[k]]$ και
$[B[1]..B[k/2]]$.
\end{enumerate}
\section{Συλλογή Comics}
\section{Πολυκατοικίες χωρίς Θέα}

\end{document}
